\chapter*{แผนบริหารการสอนประจำวิชา}
\addcontentsline{toc}{chapter}{แผนบริหารการสอนประจำวิชา}

\section*{ข้อมูลทั่วไปของรายวิชา}
\begin{itemize}[leftmargin=1.5em]
    \item \textbf{รหัสและชื่อรายวิชา} ฟิสิกส์ 2027 ฟิสิกส์แผนใหม่เชิงคำนวณและการจำลองเสมือนจริง (Computational and Immersive Modern Physics)
    \item \textbf{จำนวนหน่วยกิต} 3 หน่วยกิต (3-0-6) [บรรยาย 3 ชั่วโมง, ศึกษาด้วยตนเอง 6 ชั่วโมงต่อสัปดาห์]
    \item \textbf{หลักสูตรและประเภทของรายวิชา} วิทยาศาสตรบัณฑิต สาขาวิชาฟิสิกส์ หมวดวิชาเฉพาะด้านบังคับ
    \item \textbf{อาจารย์ผู้รับผิดชอบรายวิชาและผู้สอน} \\
    ผู้ช่วยศาสตราจารย์ ดร.ชีวะ ทัศนา (chewa.t@rbru.ac.th)
    \item \textbf{ภาคการศึกษา / ชั้นปีที่เรียน} ภาคการศึกษาที่ 2 ชั้นปีที่ 2
\end{itemize}

\section*{คำอธิบายรายวิชา (Course Description)}
ศึกษาหลักการพื้นฐานและระเบียบวิธีทางฟิสิกส์แผนใหม่ ได้แก่ ทฤษฎีสัมพัทธภาพพิเศษ สัจพจน์ของไอน์สไตน์ การแปลงลอเรนซ์ พลศาสตร์เชิงสัมพัทธภาพ กลศาสตร์ควอนตัมยุคแรก รังสีวัตถุดำ ปรากฏการณ์โฟโตอิเล็กทริก ทวิภาวะของคลื่นและอนุภาค หลักความไม่แน่นอนของไฮเซนเบิร์ก สมการชเรอดิงเงอร์และบ่อศักย์ควอนตัม แบบจำลองมาตรฐานในฟิสิกส์อนุภาค อันตรกิริยาพื้นฐาน การประยุกต์ใช้เทคโนโลยีความจริงเสริมและความจริงขยาย (AR/XR) ในการสร้างห้องปฏิบัติการเสมือนจริง โครงสร้างนิวเคลียส การสลายกัมมันตรังสีและปฏิกิริยานิวเคลียร์ หลักการคำนวณเชิงควอนตัม คิวบิตและควอนตัมเกต ตลอดจนดาราศาสตร์ฟิสิกส์ วิวัฒนาการของดาวฤกษ์ หลุมดำ และจักรวาลวิทยาร่วมกับการคำนวณเชิงตัวเลขด้วยภาษา Python

\section*{วัตถุประสงค์การเรียนรู้ (Course Learning Outcomes: CLOs) ตามอนุกรมวิธานของบลูม}
เมื่อสิ้นสุดการเรียนการสอนรายวิชานี้ นักศึกษาสามารถ
\begin{enumerate}
    \item \textbf{CLO 1 [ด้านพุทธิพิสัย: ความเข้าใจ/ความจำ]} อธิบายหลักการ ปรากฏการณ์ และสมการพื้นฐานของทฤษฎีสัมพัทธภาพ กลศาสตร์ควอนตัม ฟิสิกส์นิวเคลียร์ และจักรวาลวิทยาได้อย่างถูกต้อง
    \item \textbf{CLO 2 [ด้านพุทธิพิสัย: การประยุกต์ใช้/การวิเคราะห์]} คำนวณปริมาณทางกายภาพเชิงสัมพัทธภาพและแก้สมการชเรอดิงเงอร์สำหรับระบบศักย์อย่างง่าย พร้อมทั้งวิเคราะห์พฤติกรรมของอนุภาคได้อย่างเป็นระบบ
    \item \textbf{CLO 3 [ด้านทักษะพิสัย: การสร้างสรรค์/การคำนวณ]} พัฒนาโค้ดภาษา Python (NumPy, SciPy) เพื่อจำลองการเคลื่อนที่สัมพัทธภาพ การแก้สมการคลื่นควอนตัม และการวิวัฒนาการของระบบดาราศาสตร์ฟิสิกส์ได้
    \item \textbf{CLO 4 [ด้านทักษะพิสัย: การบูรณาการเทคโนโลยี]} ประยุกต์ใช้เฟรมเวิร์ก WebXR และ Three.js ในการสร้างและออกแบบสื่อการเรียนรู้เชิงเสมือนจริง 3 มิติ (AR/XR) ทางฟิสิกส์ได้อย่างมีประสิทธิภาพ
\end{enumerate}

\section*{แผนการจัดการเรียนรู้ตลอดภาคการศึกษา}
\begin{table}[H]
\centering
\caption*{\textbf{ตาราง}\quad แผนการจัดการเรียนรู้และการกระจายเนื้อหารายวิชาฟิสิกส์แผนใหม่}
\label{tab:modern_syllabus_map}
\small
\begin{tabularx}{\textwidth}{clp{7.5cm}c}
\toprule
\textbf{สัปดาห์} & \textbf{บทเรียน} & \textbf{หัวข้อสำคัญและกิจกรรมการเรียนรู้} & \textbf{สอดคล้อง CLO} \\
\midrule
1--2 & บทที่ 1 & บทนำ: สถาปัตยกรรมการจำลองเสมือนจริง AR/XR และการคำนวณด้วย Python & CLO 3, 4 \\
3--4 & บทที่ 2 & สัมพัทธภาพพิเศษ, การแปลงลอเรนซ์, และกรวยแสง 4D & CLO 1, 2, 3 \\
5--6 & บทที่ 3 & กลศาสตร์ควอนตัม, แบบจำลองอะตอม, สมการชเรอดิงเงอร์ และบ่อศักย์ & CLO 1, 2, 3, 4 \\
7--8 & บทที่ 4 & ฟิสิกส์ของอนุภาค, แบบจำลองมาตรฐาน, และแผนภาพไฟน์แมน & CLO 1, 2, 4 \\
9--10 & บทที่ 5 & ฟิสิกส์นิวเคลียร์, การสลายตัว, ฟิชชัน และฟิวชัน & CLO 1, 2, 3, 4 \\
11--12 & บทที่ 6 & การคำนวณเชิงควอนตัม, ทรงกลมบล็อค, และควอนตัมเกต & CLO 1, 2, 3, 4 \\
13--15 & บทที่ 7 & ดาราศาสตร์ฟิสิกส์, หลุมดำ, คลื่นความโน้มถ่วง และจักรวาลวิทยา & CLO 1, 2, 3, 4 \\
\bottomrule
\end{tabularx}
\end{table}

\section*{การวัดและประเมินผลการเรียนรู้}
\begin{itemize}[leftmargin=1.5em]
    \item การประเมินระหว่างภาคเรียน ร้อยละ 60
    \begin{itemize}
        \item แบบฝึกหัดและรายงานการวิเคราะห์เชิงตัวเลขด้วย Python ร้อยละ 20
        \item โครงงานพัฒนาสื่อจำลองเสมือนจริง AR/XR ทางฟิสิกส์ ร้อยละ 25
        \item การสอบวัดผลกลางภาคเรียน ร้อยละ 15
    \end{itemize}
    \item การสอบวัดผลปลายภาคเรียน ร้อยละ 40
\end{itemize}
